\documentclass[10pt,a4paper]{amsart}
\usepackage{amsmath,amssymb,amsthm,graphicx}
\usepackage{verbatim}
\newenvironment{reference}{\comment}{\endcomment}
%\newenvironment{reference}{}{}
\newenvironment{slogan}{\comment}{\endcomment}
\newenvironment{history}{\comment}{\endcomment}

% For commutative diagrams we use Xy-pic
\usepackage[all]{xy}

% We use 2cell for 2-commutative diagrams.
\xyoption{2cell}
% Theorem environments.
%
\theoremstyle{plain}
\newtheorem{theorem}[subsection]{Theorem}
\newtheorem{proposition}[subsection]{Proposition}
\newtheorem{lemma}[subsection]{Lemma}

\theoremstyle{definition}
\newtheorem{definition}[subsection]{Definition}
\newtheorem{example}[subsection]{Example}
\newtheorem{exercise}[subsection]{Exercise}
\newtheorem{situation}[subsection]{Situation}

\theoremstyle{remark}
\newtheorem{remark}[subsection]{Remark}
\newtheorem{remarks}[subsection]{Remarks}

\numberwithin{equation}{subsection}

% Macros
%
\def\lim{\mathop{\mathrm{lim}}\nolimits}
\def\colim{\mathop{\mathrm{colim}}\nolimits}
\def\Spec{\mathop{\mathrm{Spec}}}
\def\Hom{\mathop{\mathrm{Hom}}\nolimits}
\def\Ext{\mathop{\mathrm{Ext}}\nolimits}
\def\SheafHom{\mathop{\mathcal{H}\!\mathit{om}}\nolimits}
\def\SheafExt{\mathop{\mathcal{E}\!\mathit{xt}}\nolimits}
\def\Sch{\mathit{Sch}}
\def\Mor{\mathop{\mathrm{Mor}}\nolimits}
\def\Ob{\mathop{\mathrm{Ob}}\nolimits}
\def\Sh{\mathop{\mathit{Sh}}\nolimits}
\def\NL{\mathop{N\!L}\nolimits}
\def\CH{\mathop{\mathrm{CH}}\nolimits}
\def\proetale{{pro\text{-}\acute{e}tale}}
\def\etale{{\acute{e}tale}}
\def\QCoh{\mathit{QCoh}}
\def\Ker{\mathop{\mathrm{Ker}}}
\def\Im{\mathop{\mathrm{Im}}}
\def\Coker{\mathop{\mathrm{Coker}}}
\def\Coim{\mathop{\mathrm{Coim}}}

% Boxtimes
%
\DeclareMathSymbol{\boxtimes}{\mathbin}{AMSa}{"02}

%
% Macros for moduli stacks/spaces
%
\def\QCohstack{\mathcal{QC}\!\mathit{oh}}
\def\Cohstack{\mathcal{C}\!\mathit{oh}}
\def\Spacesstack{\mathcal{S}\!\mathit{paces}}
\def\Quotfunctor{\mathrm{Quot}}
\def\Hilbfunctor{\mathrm{Hilb}}
\def\Curvesstack{\mathcal{C}\!\mathit{urves}}
\def\Polarizedstack{\mathcal{P}\!\mathit{olarized}}
\def\Complexesstack{\mathcal{C}\!\mathit{omplexes}}
% \Pic is the operator that assigns to X its picard group, usage \Pic(X)
% \Picardstack_{X/B} denotes the Picard stack of X over B
% \Picardfunctor_{X/B} denotes the Picard functor of X over B
\def\Pic{\mathop{\mathrm{Pic}}\nolimits}
\def\Picardstack{\mathcal{P}\!\mathit{ic}}
\def\Picardfunctor{\mathrm{Pic}}
\def\Deformationcategory{\mathcal{D}\!\mathit{ef}}
\usepackage{cleveref}
\crefname{originallemma}{original lemma}{original lemmas}
\Crefname{originallemma}{Original lemma}{Original lemmas}
\makeatletter
\newlabel{coherent-lemma-flat-base-change-cohomology}{{02KH}{source}{}{sourcepointer-1}{}}
\newlabel{coherent-lemma-flat-base-change-cohomology@cref}{{[originallemma][0][]02KH}{[1][0][]source}{}{}{}}
\makeatother
\crefname{originalreference}{original reference}{original references}
\Crefname{originalreference}{Original reference}{Original references}
\makeatletter
\newlabel{example-rational-bridge}{{0E3M}{source}{}{sourcepointer-2}{}}
\newlabel{example-rational-bridge@cref}{{[originalreference][0][]0E3M}{[1][0][]source}{}{}{}}
\makeatother
\crefname{originalreference}{original reference}{original references}
\Crefname{originalreference}{Original reference}{Original references}
\makeatletter
\newlabel{lemma-no-rational-bridge-ample-genus-geq-2}{{0E3P}{source}{}{sourcepointer-3}{}}
\newlabel{lemma-no-rational-bridge-ample-genus-geq-2@cref}{{[originalreference][0][]0E3P}{[1][0][]source}{}{}{}}
\makeatother
\crefname{originalreference}{original reference}{original references}
\Crefname{originalreference}{Original reference}{Original references}
\makeatletter
\newlabel{lemma-rational-bridge-canonical}{{0E3N}{source}{}{sourcepointer-4}{}}
\newlabel{lemma-rational-bridge-canonical@cref}{{[originalreference][0][]0E3N}{[1][0][]source}{}{}{}}
\makeatother
\crefname{originalreference}{original reference}{original references}
\Crefname{originalreference}{Original reference}{Original references}
\makeatletter
\newlabel{lemma-sanity-check-duality}{{0E32}{source}{}{sourcepointer-5}{}}
\newlabel{lemma-sanity-check-duality@cref}{{[originalreference][0][]0E32}{[1][0][]source}{}{}{}}
\makeatother
\crefname{originallemma}{original lemma}{original lemmas}
\Crefname{originallemma}{Original lemma}{Original lemmas}
\makeatletter
\newlabel{morphisms-lemma-proper-ample-is-proj}{{0C6J}{source}{}{sourcepointer-6}{}}
\newlabel{morphisms-lemma-proper-ample-is-proj@cref}{{[originallemma][0][]0C6J}{[1][0][]source}{}{}{}}
\makeatother
\title{Contracting rational bridges gives a canonical model with ample canonical sheaf}
\author{The Stacks Project Authors}
\date{stacks-project:curves.tex:local-snapshot-20261012; GNU FDL 1.2 or later}
\begin{document}
\maketitle
\noindent\textit{Editorial note.} All delivered mathematical text is the author's own, copied from the pinned Stacks Project chapter file sources/stacks-curves/pinned-original-source.tex; no statement or proof step was rewritten, reordered or removed.\par\medskip
\noindent\textit{Editorial note.} The counted proof is the author's complete argument: the base-change paragraph (source lines 6526-6537) presents its argument (pullback of omega plus Gamma \textbackslash{}otimes K = Gamma\_K by flat base change, combined with the Proj canonicity argument above) and ends with 'Some details omitted.'; the ingredient f\_*O\_X = O\_\{X'\} for a single contraction is delegated to the bridge-contraction statement (source lines 6514-6516); both pointers resolve inside the chapter and the two cited lemmas are registered inputs.\par\medskip
\noindent\textit{Editorial note.} The label lemma-contracting-rational-bridges names the counted statement, and the proof environment immediately following it in the source is delivered verbatim as its proof.\par\medskip
\noindent\textit{Editorial note.} The counted statement is carried by the author's own lemma environment in the chapter source.\par\medskip
\noindent\textit{Allowed standard prerequisites (editorial).}\begin{itemize}
\item No rational bridges means ample canonical sheaf. Same chapter, \textbackslash{}label\{lemma-no-rational-bridge-ample-genus-geq-2\}. Proof lines 6508-6510: after contracting rational bridges until none are left, omega\_\{X'\} is ample.
\item Canonical pullback under a bridge contraction. Same chapter, \textbackslash{}label\{lemma-rational-bridge-canonical\}. Proof lines 6512-6514: f\textasciicircum{}*omega\_\{X'\} = omega\_X follows from this lemma together with induction over the sequence of contractions.
\item Structure sheaf pushes down under a single bridge contraction. Same chapter, Example \textbackslash{}label\{example-rational-bridge\}. Proof lines 6514-6516: f\_*O\_X = O\_\{X'\} because this is true for the contraction of a rational bridge (author's delegation).
\item Proper ample schemes are Proj of the section ring. Morphisms, Lemma \textbackslash{}label\{morphisms-lemma-proper-ample-is-proj\}. Proof lines 6520-6524: X' is the Proj of the direct sum of the spaces Gamma(X', omega\_\{X'\}\textasciicircum{}\{\textbackslash{}otimes m\}), which makes the morphism X -> X' completely canonical.
\item Dualizing module and cohomology under base field extension. Same chapter, \textbackslash{}label\{lemma-sanity-check-duality\}; Cohomology of Schemes, Lemma \textbackslash{}label\{coherent-lemma-flat-base-change-cohomology\}. Proof lines 6526-6537: omega\_\{X\_K\} is the pullback of omega\_X and Gamma(X, omega\_X\textasciicircum{}\{\textbackslash{}otimes m\}) \textbackslash{}otimes\_k K = Gamma(X\_K, omega\_\{X\_K\}\textasciicircum{}\{\textbackslash{}otimes m\}) by flat base change.
\end{itemize}\medskip
\ifdefined\hypertarget\hypertarget{sourcepointer-1}{}\fi
\noindent\textit{Source reference.} Original lemma 02KH: sources/stacks-chow/tags.txt, source lines 3372--3372. Prerequisite pointer invoked at proof line(s) 6536; it names a standard fact of a sibling chapter that the counted argument uses, and every step of the counted proof built on it is retained verbatim in the delivered body.\par\medskip
\ifdefined\hypertarget\hypertarget{sourcepointer-2}{}\fi
\noindent\textit{Source reference.} Original reference 0E3M: sources/stacks-curves/pinned-original-source.tex, source lines 6201--6201. Prerequisite pointer invoked at proof line(s) 6500; it names a standard fact of the same chapter that the counted argument uses, and every step of the counted proof built on it is retained verbatim in the delivered body.\par\medskip
\ifdefined\hypertarget\hypertarget{sourcepointer-3}{}\fi
\noindent\textit{Source reference.} Original reference 0E3P: sources/stacks-curves/pinned-original-source.tex, source lines 6433--6433. Prerequisite pointer invoked at proof line(s) 6509; it names a standard fact of the same chapter that the counted argument uses, and every step of the counted proof built on it is retained verbatim in the delivered body.\par\medskip
\ifdefined\hypertarget\hypertarget{sourcepointer-4}{}\fi
\noindent\textit{Source reference.} Original reference 0E3N: sources/stacks-curves/pinned-original-source.tex, source lines 6367--6367. Prerequisite pointer invoked at proof line(s) 6513; it names a standard fact of the same chapter that the counted argument uses, and every step of the counted proof built on it is retained verbatim in the delivered body.\par\medskip
\ifdefined\hypertarget\hypertarget{sourcepointer-5}{}\fi
\noindent\textit{Source reference.} Original reference 0E32: sources/stacks-curves/pinned-original-source.tex, source lines 927--927. Prerequisite pointer invoked at proof line(s) 6531; it names a standard fact of the same chapter that the counted argument uses, and every step of the counted proof built on it is retained verbatim in the delivered body.\par\medskip
\ifdefined\hypertarget\hypertarget{sourcepointer-6}{}\fi
\noindent\textit{Source reference.} Original lemma 0C6J: sources/stacks-chow/tags.txt, source lines 15299--15299. Prerequisite pointer invoked at proof line(s) 6525; it names a standard fact of a sibling chapter that the counted argument uses, and every step of the counted proof built on it is retained verbatim in the delivered body.\par\medskip
\begin{lemma}
\label{lemma-contracting-rational-bridges}
Let $k$ be a field. Let $X$ be a proper scheme over $k$ of dimension $1$
with $H^0(X, \mathcal{O}_X) = k$ having genus $g \geq 2$.
Assume the singularities of $X$ are at-worst-nodal and that
$X$ has no rational tails. Consider a sequence
$$
X = X_0 \to X_1 \to \ldots \to X_n = X'
$$
of contractions of rational bridges
(Example \ref{example-rational-bridge}) until none are left.
Then $\omega_{X'}$ ample.
The morphism $X \to X'$ is independent of choices and
formation of this morphism commutes with base field extensions.
\end{lemma}

\begin{proof}
We proceed by contracting rational bridges until there are none
left. Then $\omega_{X'}$ is ample by
Lemma \ref{lemma-no-rational-bridge-ample-genus-geq-2}.

\medskip\noindent
Denote $f : X \to X'$ the composition. By
Lemma \ref{lemma-rational-bridge-canonical} and induction we see that
$f^*\omega_{X'} = \omega_X$.
We have $f_*\mathcal{O}_X = \mathcal{O}_{X'}$
because this is true for contraction of a rational bridge.
Thus the projection formula says that
$f_*f^*\mathcal{L} = \mathcal{L}$ for all invertible
$\mathcal{O}_{X'}$-modules $\mathcal{L}$.
Hence
$$
\Gamma(X', \omega_{X'}^{\otimes m}) = \Gamma(X, \omega_X^{\otimes m})
$$
for all $m$. Since $X'$ is the Proj of the direct sum of these
by Morphisms, Lemma \ref{morphisms-lemma-proper-ample-is-proj}
we conclude that the morphism $X \to X'$ is completely canonical.

\medskip\noindent
Let $K/k$ be an extension of fields, then
$\omega_{X_K}$ is the pullback of $\omega_X$
(Lemma \ref{lemma-sanity-check-duality}) and we have
$\Gamma(X, \omega_X^{\otimes m}) \otimes_k K$
is equal to
$\Gamma(X_K, \omega_{X_K}^{\otimes m})$
by
Cohomology of Schemes, Lemma \ref{coherent-lemma-flat-base-change-cohomology}.
Thus formation of $f : X \to X'$ commutes with base change by
$K/k$ by the arguments given above. Some details omitted.
\end{proof}

\end{document}
